\scanpage{84}\mspage{41}
priori que cette condition supplémentaire
est inutile \deleted{\unsure{q}}pour impliquer que $\widehat u$ est
$\Lambda$-linéaire, quand $\Lambda$ est un anneau
absolu $(\simeq\mathbb Z\ \text{ou}\ \mathbb Z/n\mathbb Z)$, puisque tt
autom. additif d'un $\Lambda$-module est
alors au\deleted{\unsure{toma|tomat}}tomatiquement $\Lambda$-linéaire. Par
contre on comprend bien pourquoi,
quand $\Lambda$ n'est pas un anneau absolu,
il faut absolument mettre une structure
supplémentaire sur $\widetilde P_M$ \inserted{$=(\unsure{S_M|\widetilde S_M},\widetilde F_M)$} (genre op. de $\Lambda^*$
dessus) -- en effet tout autom. \uline{additif}
(pas néc. $\Lambda$-linéaire) de $M$ induit un
autom. de $\widetilde P_M$ ! Mais cette remarque
ne s'applique pas telle quelle, si on
tient compte de la structure suppl.
$\widetilde F_M=\coprod_i\widetilde F_{M,i}$ -- \deleted{sauf si} \inserted{si pourtant \unsure{si|si si}} $\Lambda^*=\{\pm1\}$
(p.~ex. $\Lambda=$ \deleted{$\mathbb F_3$} anneau de polynômes sur
$\mathbb F_2,\mathbb F_3$ ou $\mathbb Z$), auquel cas l'ens. d'indices
$\pi$ lui-m$^{\text{me}}$ est réduit à un élément !
% archive sheet number: 83
